% GENERATED FILE. Edit canonical lessons, then run npm run generate:books.
% canonical-source-hash: fnv1a64:72b9e50a4f2aea01
% canonical-lessons: GU-C384-saphal, GU-C384-unnati, GU-C384-vikas, GU-C384-phavvu, GU-C384-pahocadvu

\chapter{Success}
\label{ch:gu-success}
\hlchaptermodality{384}

\begin{chapteropening}
\textbf{By the end of this chapter:} \emph{I can say what it will take to succeed.}

Complete the last lesson of chapter 384: I can say what it will take to succeed.
\end{chapteropening}

\section[\texorpdfstring{\gu{સફળ} (saphaḷ)}{સફળ (saphaḷ)}]{\gu{સફળ} (saphaḷ) — successful}

\label{lesson:GU-C384-saphal}



\begin{quote}
Before the new one: say the Gujarati for to set about, then the Gujarati for to battle on.
\end{quote}

\subsection*{You'll want to know: \gu{સફળ}}
\textbf{\gu{સફળ}} — \emph{saphaḷ} — \textquotedblleft{}successful\textquotedblright{}.

\begin{quote}
\gu{જો} \gu{આ} \gu{યોજના} \gu{સફળ} \gu{થશે} \gu{તો} \gu{આખા} \gu{ગામને} \gu{ફાયદો} \gu{થશે}.
\end{quote}

\begin{quote}
\emph{jo ā yojnā saphaḷ thaśe to ākhā gāmne phāydo thaśe.}
\end{quote}

\begin{quote}
If this plan succeeds, the whole village will benefit.
\end{quote}

\begin{cousinweb}
\textbf{\gu{સફળ}} (\emph{saphaḷ}) is a tatsama, borrowed from Sanskrit \emph{saphala}, \textquotedblleft{}having fruit, fruitful\textquotedblright{} (Monier-Williams).
\end{cousinweb}

\begin{grammarlens}[title={What you've built}]
One more word for saying what happens if.
\end{grammarlens}

\subsection*{Your turn}
\begin{itemize}
\raggedright
  \item \emph{Say it:} \emph{saphaḷ}
  \item \emph{Say it:} \emph{saphaḷ}, once more
  \item \emph{Say it:} \emph{jo ā yojnā saphaḷ thaśe to ākhā gāmne phāydo thaśe.} — If this plan succeeds, the whole village will benefit.
  \item \emph{Recall:} say the Gujarati for to set about, then the Gujarati for to battle on, then say \emph{saphaḷ} again
\end{itemize}

\begin{tcolorbox}[breakable,colback=teal!4,colframe=teal!35!black,title={Before you move on}]
What does \gu{સફળ} mean? (\textquotedblleft{}Successful\textquotedblright{}.) Say it once more.
\end{tcolorbox}
\section[\texorpdfstring{\gu{ઉન્નતિ} (unnati)}{ઉન્નતિ (unnati)}]{\gu{ઉન્નતિ} (unnati) — advancement}

\label{lesson:GU-C384-unnati}



\begin{quote}
Before the new one: say the Gujarati for to battle on, then the Gujarati for successful.
\end{quote}

\subsection*{You'll want to know: \gu{ઉન્નતિ}}
\textbf{\gu{ઉન્નતિ}} — \emph{unnati} — \textquotedblleft{}advancement\textquotedblright{}.

\begin{quote}
\gu{જો} \gu{ગામમાં} \gu{શાળા} \gu{અને} \gu{દવાખાનું} \gu{હશે} \gu{તો} \gu{ગામની} \gu{ઉન્નતિ} \gu{થશે}.
\end{quote}

\begin{quote}
\emph{jo gāmmā̃ śāḷā ane davākhānũ haśe to gāmnī unnati thaśe.}
\end{quote}

\begin{quote}
If the village has a school and a clinic, it will move ahead.
\end{quote}

A feminine noun.

\begin{cousinweb}
\textbf{\gu{ઉન્નતિ}} (\emph{unnati}) is a tatsama, borrowed from Sanskrit \emph{unnati}, \textquotedblleft{}rising, ascending\textquotedblright{} (Monier-Williams).
\end{cousinweb}

\begin{grammarlens}[title={What you've built}]
One more word for saying what happens if.
\end{grammarlens}

\subsection*{Your turn}
\begin{itemize}
\raggedright
  \item \emph{Say it:} \emph{unnati}
  \item \emph{Say it:} \emph{unnati}, once more
  \item \emph{Say it:} \emph{jo gāmmā̃ śāḷā ane davākhānũ haśe to gāmnī unnati thaśe.} — If the village has a school and a clinic, it will move ahead.
  \item \emph{Recall:} say the Gujarati for to battle on, then the Gujarati for successful, then say \emph{unnati} again
\end{itemize}

\begin{tcolorbox}[breakable,colback=teal!4,colframe=teal!35!black,title={Before you move on}]
What does \gu{ઉન્નતિ} mean? (\textquotedblleft{}Advancement\textquotedblright{}.) Say it once more.
\end{tcolorbox}
\section[\texorpdfstring{\gu{વિકાસ} (vikās)}{વિકાસ (vikās)}]{\gu{વિકાસ} (vikās) — development}

\label{lesson:GU-C384-vikas}



\begin{quote}
Before the new one: say the Gujarati for successful, then the Gujarati for advancement.
\end{quote}

\subsection*{You'll want to know: \gu{વિકાસ}}
\textbf{\gu{વિકાસ}} — \emph{vikās} — \textquotedblleft{}development\textquotedblright{}.

\begin{quote}
\gu{જો} \gu{સારા} \gu{રસ્તા} \gu{બનશે} \gu{તો} \gu{શહેરનો} \gu{વિકાસ} \gu{ઝડપી} \gu{થશે}.
\end{quote}

\begin{quote}
\emph{jo sārā rastā banśe to śaherno vikās jhaḍpī thaśe.}
\end{quote}

\begin{quote}
If good roads are built, the city will develop faster.
\end{quote}

A masculine noun.

\begin{cousinweb}
\textbf{\gu{વિકાસ}} (\emph{vikās}) is a tatsama, borrowed from Sanskrit \emph{vikāsa}, \textquotedblleft{}expanding, blossoming\textquotedblright{} (Monier-Williams).
\end{cousinweb}

\begin{grammarlens}[title={What you've built}]
One more word for saying what happens if.
\end{grammarlens}

\subsection*{Your turn}
\begin{itemize}
\raggedright
  \item \emph{Say it:} \emph{vikās}
  \item \emph{Say it:} \emph{vikās}, once more
  \item \emph{Say it:} \emph{jo sārā rastā banśe to śaherno vikās jhaḍpī thaśe.} — If good roads are built, the city will develop faster.
  \item \emph{Recall:} say the Gujarati for successful, then the Gujarati for advancement, then say \emph{vikās} again
\end{itemize}

\begin{tcolorbox}[breakable,colback=teal!4,colframe=teal!35!black,title={Before you move on}]
What does \gu{વિકાસ} mean? (\textquotedblleft{}Development\textquotedblright{}.) Say it once more.
\end{tcolorbox}
\section[\texorpdfstring{\gu{ફાવવું} (phāvvũ)}{ફાવવું (phāvvũ)}]{\gu{ફાવવું} (phāvvũ) — to get one's way; to succeed}

\label{lesson:GU-C384-phavvu}



\begin{quote}
Before the new one: say the Gujarati for advancement, then the Gujarati for development.
\end{quote}

\subsection*{You'll want to know: \gu{ફાવવું}}
\textbf{\gu{ફાવવું}} — \emph{phāvvũ} — \textquotedblleft{}to get one's way; to succeed\textquotedblright{}.

\begin{quote}
\gu{જો} \gu{તે} \gu{છેતરવાની} \gu{કોશિશ} \gu{કરશે} \gu{તો} \gu{તે} \gu{નહીં} \gu{ફાવે}.
\end{quote}

\begin{quote}
\emph{jo te chetarvānī kośiś karśe to te nahī̃ phāve.}
\end{quote}

\begin{quote}
If he tries to cheat, he will not get away with it.
\end{quote}

Past (he …): \gu{તે} \gu{પોતાની} \gu{યોજનામાં} \gu{ફાવ્યો}, \emph{te potānī yojnāmā̃ phāvyo}.

\begin{cousinweb}
\textbf{\gu{ફાવવું}} (\emph{phāvvũ}) CDIAL lists the Gujarati verb, \textquotedblleft{}to find convenient, to succeed\textquotedblright{}, under a reconstructed \emph{sparvati}, \textquotedblleft{}releases, wins, succeeds\textquotedblright{} (CDIAL 13808).
\end{cousinweb}

\begin{grammarlens}[title={What you've built}]
One more word for saying what happens if.
\end{grammarlens}

\subsection*{Your turn}
\begin{itemize}
\raggedright
  \item \emph{Say it:} \emph{phāvvũ}
  \item \emph{Say it:} \emph{phāvvũ}, once more
  \item \emph{Say it:} \emph{jo te chetarvānī kośiś karśe to te nahī̃ phāve.} — If he tries to cheat, he will not get away with it.
  \item \emph{Recall:} say the Gujarati for advancement, then the Gujarati for development, then say \emph{phāvvũ} again
\end{itemize}

\begin{tcolorbox}[breakable,colback=teal!4,colframe=teal!35!black,title={Before you move on}]
What does \gu{ફાવવું} mean? (\textquotedblleft{}To get one's way; to succeed\textquotedblright{}.) Say it once more.
\end{tcolorbox}
\section[\texorpdfstring{\gu{પહોંચાડવું} (pahõcāḍvũ)}{પહોંચાડવું (pahõcāḍvũ)}]{\gu{પહોંચાડવું} (pahõcāḍvũ) — to deliver}

\label{lesson:GU-C384-pahocadvu}



\begin{quote}
Before the new one: say the Gujarati for development, then the Gujarati for to get one's way.
\end{quote}

\subsection*{You'll want to know: \gu{પહોંચાડવું}}
\textbf{\gu{પહોંચાડવું}} — \emph{pahõcāḍvũ} — \textquotedblleft{}to deliver\textquotedblright{}.

\begin{quote}
\gu{જો} \gu{ટપાલી} \gu{સમયસર} \gu{આવશે} \gu{તો} \gu{તે} \gu{પત્ર} \gu{આજે} \gu{પહોંચાડશે}.
\end{quote}

\begin{quote}
\emph{jo ṭapālī samaysar āvśe to te patra āje pahõcāḍśe.}
\end{quote}

\begin{quote}
If the postman comes on time, he will deliver the letter today.
\end{quote}

Past, with \gu{તેણે} (he or she …): \gu{તેણે} \gu{પાર્સલ} \gu{ઘરે} \gu{પહોંચાડ્યું}, \emph{teṇe pārsal ghare pahõcāḍyũ}.

\begin{cousinweb}
\textbf{\gu{પહોંચાડવું}} (\emph{pahõcāḍvũ}) is the causative of \textbf{\gu{પહોંચવું}} (\emph{pahõcvũ}), \textquotedblleft{}to arrive\textquotedblright{}, which CDIAL lists under Sanskrit \emph{prábhūta}, \textquotedblleft{}much, great\textquotedblright{} (CDIAL 8716): to make a thing arrive.
\end{cousinweb}

\begin{grammarlens}[title={What you've built}]
One more word for saying what happens if.
\end{grammarlens}

\subsection*{Your turn}
\begin{itemize}
\raggedright
  \item \emph{Say it:} \emph{pahõcāḍvũ}
  \item \emph{Say it:} \emph{pahõcāḍvũ}, once more
  \item \emph{Say it:} \emph{jo ṭapālī samaysar āvśe to te patra āje pahõcāḍśe.} — If the postman comes on time, he will deliver the letter today.
  \item \emph{Recall:} say the Gujarati for development, then the Gujarati for to get one's way, then say \emph{pahõcāḍvũ} again
\end{itemize}

\begin{tcolorbox}[breakable,colback=teal!4,colframe=teal!35!black,title={Before you move on}]
What does \gu{પહોંચાડવું} mean? (\textquotedblleft{}To deliver\textquotedblright{}.) Say it once more.
\end{tcolorbox}
